% Desenvolvido por André Leite - DE - UFPE
% Ajustes Raydonal Ospina - DE - UFPE/UFBA

\documentclass[a4paper,12pt,oneside,twocolumn]{Config/milktest}

%---------------------- customização ----------------------



\usepackage{fontspec}
\usepackage{xunicode}
\usepackage{xltxtra}
\usepackage[LGR,T1]{fontenc}
\usepackage{ebgaramond}

%\usepackage[cp1252]{inputenc}                %% Use um dos dois esquemas
\usepackage[english,brazil]{babel}                 %% Hifenizacao e estilo
\usepackage{pstricks}                          %% Pacote de comandos p/ Facilitar o Postscripts
\usepackage{pstricks-add}                    %% " "
\usepackage{pst-all,pst-math,pst-plot, pst-func}       %% " "
\usepackage{textcomp}							%%
%\usepackage[tight,nice]{units}				%%
\usepackage[centertags]{amsmath}       %% [leqno]
\usepackage{amsfonts}                         %% Pacotes da AMS
\usepackage{amssymb}                         %%
\usepackage{stmaryrd} 							%%
\usepackage{Config/milkstyle}
%\usepackage{sabon}

\usepackage{nicefrac}%  Nice Frescura
\usepackage{balance}%
\usepackage{lastpage} %                      %% Referênciar a última página
\usepackage{graphicx}                        %% Mais ferramentas para
%\usepackage[T1]{pbsi} % Procurar e instalar no Mac!!!!! <------<-------<------
\usepackage{mathrsfs}                        %% $\mathscr{ABCD}$
\usepackage{lettrine}                        %% ...
%\usepackage[normal,bf,up]{caption2}          %% Controlar melhor os captions....
\usepackage{eucal}                           %% beau
\usepackage{colortbl} %
\usepackage{hyperref}                 %%
%\usepackage{amsthm}
\usepackage{indentfirst}%
\usepackage{enumitem}%
\usepackage{colordvi}                        %% Usar Cores...
\usepackage{fancyhdr}%
\usepackage{yfonts}%
\usepackage{theorem}%
\usepackage{multirow}                        %% Juntar linhas em tabelas
\usepackage{midfloat}%
\usepackage{cancel}%
%\usepackage[square]{natbib}%sectionbib,
\usepackage{fancyhdr} % cabeçalhos ...
\usepackage{rotating}


\usepackage{lipsum} 
%% Lista de Pacotes utilizados


%% Deixando o caption mais interessante...
%\renewcommand{\captionfont}{\itshape}


\hypersetup{%backref,%
%            pdfpagemode=FullScreen,%
            colorlinks=true,%
            citecolor=blue,%
            urlcolor=blue, %
            % urlcolor=black, %
            linkcolor=blue,%
            pdfauthor={André Leite Wanderley}, %
            pdftitle={1EE/ET406}, %
            pdfcreator={André Leite Wanderley}, %
            breaklinks=true, %
            linktocpage=true
            }%


\DeclareMathOperator*{\Inf}{Inf}%
\DeclareMathOperator*{\Sup}{Sup}%
\DeclareMathOperator*{\Min}{Min}%
\DeclareMathOperator*{\Max}{Max}
\DeclareMathOperator*{\Sgn}{Sgn}%
\DeclareMathOperator*{\E}{E}%
\newcommand{\mbf}{\boldsymbol}%
\DeclareMathAlphabet\mathbb{U}{fplmbb}{m}{n}%
%
%\numberwithin{equation}{section}                       %%
%\numberwithin{table}{section}                          %% Numeração dos floats...
%\numberwithin{figure}{section}                         %%


\newcommand{\oo}{\mathbin{\bullet}}
\newcommand{\oc}{\mathbin{\circ}}
\newcommand{\noo}{\mathbin{\not\!\bullet}}
\newcommand{\cf}{\emph{cf.}\ }
\newcommand{\ie}{\emph{i.e.}}
\newcommand{\eg}{\emph{e.g.}}
\newcommand{\raab}{\emph{reductio ad absurdum}}
\newcommand{\eetc}{\textit{et c\ae tera}}
\newcommand{\dom}{\mathfrak{D}}
\newcommand{\ima}{\mathfrak{I}}
\newcommand{\idenA}{\mathcal{I}_\mathcal{A}}
\newcommand{\ce}{\ensuremath{\mathcal{X}}}
\newcommand{\piq}{\ensuremath{\succsim}}
\newcommand{\pq}{\ensuremath{\succ}}
\newcommand{\iq}{\ensuremath{\sim}}
\newcommand{\niq}{\ensuremath{\nsim}}
\newcommand{\npiq}{\ensuremath{\not\succsim}}
\newcommand{\npq}{\ensuremath{\nsucc}}
\newcommand{\po}[2]{\ensuremath{(#1,#2)}}

%\setlist{noitemsep}
%\setlist{labelindent=\parindent} % << Usually a good idea
\setlist{itemsep=0pt}%
\setitemize{leftmargin=\parindent} %
\setitemize[1]{label=$\triangleright$} %
%\setenumerate[1]{label=\arabic*.,ref=\arabic*}
\setenumerate[1]{label=\arabic*.,ref=\arabic*}
\setenumerate[2]{label=\theenumi.\arabic*.,ref=\theenumi.\arabic*}
\setenumerate[3]{label=\roman*.,ref=\theenumii.\roman*}
%\setdescription{%font=\sffamily\bfseries}

%                }%


%%%%%%%%%%%%%%%%%%%%%%%%%%%%%%%%
% Para facilitar minha vida

\newcommand{\Ds}{\displaystyle}%
\newcommand{\Ts}{\textstyle}%
\newcommand{\Ss}{\scriptstyle}%
\newcommand{\SSs}{\scriptscriptstyle}%

\def\aureos{%
            \newrgbcolor{gray1}{.9 .9 .9}
            \newrgbcolor{guess1}{1 0.388235294 0.090196078}
            \psset{unit=1in,cornersize=absolute}%
            \psline[linecolor=guess1,linewidth=1.2pt](-0.125,-1.291667)(1.1125,-1.291667)
            \psline[linecolor=guess1,linewidth=0.6pt](0.347683,-1.291667)(0.347683,-0.52685)
            \psline[linecolor=guess1,linewidth=0.6pt](0.347683,-0.999533)(-0.125,-0.999533)
            \psline[linecolor=guess1,linewidth=0.6pt](0.167134,-0.999533)(0.167134,-1.291667)
            \psline[linecolor=guess1,linewidth=0.6pt](0.167134,-1.111118)(0.347683,-1.111118)
            \psline[linecolor=guess1,linewidth=0.6pt](0.236098,-1.111118)(0.236098,-0.999533)
            \psline[linecolor=guess1,linewidth=0.6pt](0.236098,-1.068496)(0.167134,-1.068496)
            \psline[linecolor=guess1,linewidth=0.6pt](0.209756,-1.068496)(0.209756,-1.111118)
            \psline[linecolor=guess1,linewidth=0.6pt](0.209756,-1.084776)(0.236098,-1.084776)
           }%

\def\GuEsS{%
            \newrgbcolor{guess1}{1 0.388235294 0.090196078}
            \newrgbcolor{guess2}{0.960784314 0.109803922 0.043137255}
            \psset{unit=1in,cornersize=absolute}%
            \rput(0,-1.5){{\fontsize{35}{0}\usefont{OT1}{cmr}{m}{n}\selectfont\textsc{\textbf{S}}}}%
            \rput(0.625,-1.5){\shortstack{{\huge\textbf{\textsc{ystems}}}\\
                {\scriptsize\sc\hfill Research Group}}}
            \psline[linecolor=black,fillcolor=guess2,fillstyle=none,linewidth=3.6pt](-0.125,-1.725)(1.1125,-1.725)%
            \uput{0.5ex}[l](1.148214,-1.229167){\llap{{\tiny\sc UFBA}}}%
            \psline[linecolor=guess1,fillcolor=guess1,fillstyle=none,linewidth=1.2pt](-0.125,-1.291667)(1.1125,-1.291667)%
            \psline[linecolor=guess1,fillcolor=guess1,fillstyle=none,linewidth=.6pt](0.347683,-1.291667)(0.347683,-0.52685)%
            \psline[linecolor=guess1,fillcolor=guess1,fillstyle=none,linewidth=.6pt](0.347683,-0.999533)(-0.125,-0.999533)%
            \psline[linecolor=guess1,fillcolor=guess1,fillstyle=none,linewidth=.6pt](0.167134,-0.999533)(0.167134,-1.291667)%
            \psline[linecolor=guess1,fillcolor=guess1,fillstyle=none,linewidth=.6pt](0.167134,-1.111118)(0.347683,-1.111118)%
            \psline[linecolor=guess1,fillcolor=guess1,fillstyle=none,linewidth=.6pt](0.236098,-1.111118)(0.236098,-0.999533)%
            \psline[linecolor=guess1,fillcolor=guess1,fillstyle=none,linewidth=.6pt](0.236098,-1.068496)(0.167134,-1.068496)%
            \psline[linecolor=guess1,fillcolor=guess1,fillstyle=none,linewidth=.6pt](0.209756,-1.068496)(0.209756,-1.111118)%
            \psline[linecolor=guess1,fillcolor=guess1,fillstyle=none,linewidth=.6pt](0.209756,-1.084776)(0.236098,-1.084776)%
           }%

%%%%%%%%%%%%%%%%%%%%%%%%%%%%%%%%%%%%%%%%%%%%%%%%%%%%%%%%


\theoremstyle{plain}                     %%
%\theoremstyle{margin}                   %%
\theoremheaderfont{\scshape\bfseries}    %%
\theorembodyfont{\rmfamily}                       %%

\newtheorem{Prop}{\small Proposition}[section]  %%
\newtheorem{Def}{\small Definition}[section]    %%
\newtheorem{Not}{\small Notice}[section]    %%
\newtheorem{Theo}{\small Theorem}[section]      %%
\newtheorem{Lem}{\small Lema}[section]         %%
\newtheorem{Coro}{\small Corollary}[section]    %%
\newtheorem{Conj}{\small Conjecture}[section]   %%
\newtheorem{Prob}{\small Problem}[section]      %%
\newtheorem{Exa}{\small Example}[section]       %%
%\theoremheaderfont{\bfseries}    %%
%\theorembodyfont{\rmfamily}  %%
\newtheorem{proof}{\small Demonstration}[section]
\newtheorem{sol}{\small Solution}[section]


\newenvironment{Proofm}               % Deve ter uma maneira
    {%\strut \par\medskip             % ridiculamente mais
     %\strut\hspace{\parindent}      % fácil de fazer isso...
     %\begin{minipage}{.8\linewidth} % I'm sure about that.
     \begin{proof}\small}                   %
    {\strut\hfill\textsc{q.e.d.}\end{proof}                      %
    %\end{minipage}                  %
    %\strut\par\medskip
                        }             %

\newenvironment{solm}               % Deve ter uma maneira
    {%\strut\par%\medskip             % ridiculamente mais
     %\strut\hspace{\parindent}      % fácil de fazer isso...
     %\begin{minipage}{.8\linewidth} % I'm sure about that.
     %\begin{sol}
     \small
     \noindent\textbf{Solução:}
     %\par %\smallskip
     }                   %
    {\strut\hfill\textsc{$\blacksquare$}%\end{sol}                      %
    %\end{minipage}                  %
    %\strut\par\medskip
                        }             %



\newenvironment{Propm}               % Deve ter uma maneira
    {\strut\par\smallskip             % ridiculamente mais
     \strut\hspace{\parindent}      % fácil de fazer isso...
     \begin{minipage}{.8\linewidth} % I'm sure about that.
     \begin{Prop}\small}                   %
    {\end{Prop}                      %
    \end{minipage}                  %
    \strut\par\medskip
    }             %


\newenvironment{Defm}               % Deve ter uma maneira
    {\strut\par\smallskip             % ridiculamente mais
     \strut\hspace{\parindent}      % fácil de fazer isso...
     \begin{minipage}{.8\linewidth} % I'm sure about that.
     \begin{Def}\small}                   %
    {\end{Def}                      %
    \end{minipage}                  %
    \strut\par\medskip}             %

\newenvironment{Theom}               % Deve ter uma maneira
    {\strut\par\smallskip             % ridiculamente mais
     \strut\hspace{\parindent}      % fácil de fazer isso...
     \begin{minipage}{.8\linewidth} % I'm sure about that.
     \begin{Theo}\small}                   %
    {\end{Theo}                      %
    \end{minipage}                  %
    \strut\par\medskip}

\newenvironment{Notm}               % Deve ter uma maneira
    {\strut\par\smallskip             % ridiculamente mais
     \strut\hspace{\parindent}      % fácil de fazer isso...
     \begin{minipage}{.8\linewidth} % I'm sure about that.
     \begin{Not}\small}                   %
    {\end{Not}                      %
    \end{minipage}                  %
    \strut\par\medskip}


\newenvironment{Exam}               % Deve ter uma maneira
    {%\strut\par\smallskip             % ridiculamente mais
     %\strut\hspace{\parindent}      % fácil de fazer isso...
     %\begin{minipage}{1\linewidth} % I'm sure about that.
     \begin{Exa}\small}                   %
    { \scriptsize$\blacksquare$\end{Exa}                      %
    %\end{minipage}                  %
    \strut\par\medskip}

\def\Pt#1(#2,#3){
       \expandafter\def\csname #1\endcsname{#2,#3}%
       \expandafter\def\csname #1X\endcsname{#2}%
       \expandafter\def\csname #1Y\endcsname{#3}}%

\newcommand{\hlinem}[3]{\noalign{\vspace{#3}\hrule height #1 depth #2}}
\newcommand{\hlinemilk}[3]{\noalign{}}

\newcommand{\quem}[1]{\par\medskip\strut\hfill\bf  --- #1}
\newcounter{questionNumber}
\newcommand{\question}{\stepcounter{questionNumber}\lettrine[lhang=.75, findent=10pt, nindent=0pt,lines=3, lraise=.2]{\thequestionNumber}{$\blacktriangleright$}}


\newenvironment{milkquote}
    {\medskip %
    \begin{center}
    \begin{minipage}{.9\linewidth}
    \hspace{.25cm}\small\itshape}%
    {\medskip\end{minipage}
    \end{center}
    }%


\newcommand{\fim}{%
                 \par\medskip
                 \begin{center}
                      \includegraphics[width=1cm]{ornaments/ornamento4.eps}
                 \end{center}
                 \vspace*{.5cm}
                 }%



\newenvironment{milkanswer}
    {\par %
    \begin{center}
    \color{orange}\textbullet\quad\textbf{Solution}\quad\textbullet
    \end{center}\small
    }%
    {
    \fim
    }


%%%
%%% Set up page layout parameters.
%%%
\setlength\paperheight{297mm}      %
\setlength\paperwidth{210mm}       % papel a4
%\setlength{\parindent}{6mm}

\addtolength{\hoffset}{-1in}%
\addtolength{\voffset}{-1.5in}

\setlength{\oddsidemargin}{25mm}    % 
\setlength{\topmargin}{5cm}        %
\setlength{\textwidth}{170mm}      % Margem direita: 25mm
\setlength{\textheight}{220mm}  % Margem Inferior: 25mm
\setlength{\marginparsep}{0mm}     %

\setlength{\headheight}{16mm}
\renewcommand{\columnsep}{1cm}%
%\marginsize{1.25in}{1.75in}{1.25in}{1.25in}



% If you can to leave the header here
%---

%---

\renewcommand{\headrulewidth}{0pt}
\definecolor{cor}{rgb}{.5,0, 0}

%%%%%%%%%%%%%%%%%%%%%%%%%%% Font 

\defaultfontfeatures{Ligatures=Historic,Contextuals=Alternate, Numbers=OldStyle,RawFeature={ +hlig, +dlig, +ss05}}
\setmainfont[RawFeature={+ss05, +hlig, +dlig}, BoldFont = Gentium Basic Bold, ItalicFeatures={RawFeature={+clig},CharacterVariant=5:0}]{EB Garamond}
% % Do the replacements manually in titles, to not ruin the textsc
% \newfontfamily\booktitlefont[RawFeature={-ss02},LetterSpace=40,WordSpace=6]{EB Garamond}
% \newfontfamily\spacedfont[RawFeature={-ss02},LetterSpace=20,WordSpace=3]{EB Garamond}
% \newfontfamily\lettrinefont{EB Garamond Initials}
% \newfontfamily\headerfont{EB Garamond}

%\setmainfont[BoldFont = Gentium Basic Bold,
%      ItalicFont = Gentium Plus Italic,
%      BoldItalicFont = Gentium Basic Bold Italic,
%      Mapping = tex-text]{Gentium Plus}





 

\usepackage{bm}
\usepackage{verbatimbox}
\usepackage{float}

%%%%%%%%%%%%%----- Header------%%%%%%%%%%%%%%%%%%%
\newsavebox{\mygraphic}
\sbox{\mygraphic}{\includegraphics[scale=1.2]{Figuras/ufba.eps}}

\fancypagestyle{milktablestyle}{
    \fancyhead{}
    \fancyfoot[C]{\thepage/\pageref{LastPage}} %
}

\pagestyle{fancy}

\fancyhead{}
\fancyhead[C]{
\rput(0,2){
  \begin{minipage}{1.75cm}\centering
    \includegraphics[width=1.75cm]{Figuras/ufbaBlack.eps}
  \end{minipage}~\hfill~\begin{minipage}{10cm}\Large\centering
   \color{cor}{\scshape universidade federal da bahia}\\
    {\scshape instituto de matemática e estatística}\\
    {\scshape departamento de estatística}\\[5pt]
    % Aqui as informações da prova e a data
    {\Large\scshape  prova I --- matd44 --- gabarito} 
  \end{minipage}~\hfill~\begin{minipage}{1.75cm}\centering
    \includegraphics[width=1.75cm]{Figuras/ufba.eps}
  \end{minipage}
  }}

\fancyfoot[C]{\thepage/\pageref{LastPage}} % 

%%%%%%%%%%%%%----- Header------%%%%%%%%%%%%%%%%%%%

% \usepackage[printwatermark]{xwatermark}
% \newwatermark[allpages,color=gray!20,angle=45,scale=3,xpos=30,ypos=-60]{DRAFT}

\begin{document}




% %------------------------------------------------
\beb
{\small
\begin{description}
\item[Professor:] Raydonal Ospina
% \item[E-mail:] raydonalmartinez@ufba.br
% \item[Site do curso:] \url{https://castlab.org/courses/matd49/}

% \item[Regras:] Leia com atenção as perguntas. A prova deve ser claramente resolvida. 
% Seja claro e organizado. Pode fazer uso de computadora, calculadora, celular, I.A., livros, material de aula.

% \item[Atenção:] Responda única e exclusivamente o número de questões {\bf suficientes} para completar 10 (dez) pontos.  Fique a vontade para escolher as perguntas que achar mais interessantes. Especifique sua escolha, por exemplo, 1(a). (isto quer dizer que você escolheu  o numeral (a) da questão 1).
\end{description}
}
\eeb

% %%%%%%%%%%%%%%%%%%%%%%%%%%%%%%%%%%%%%%%%%%%%%%%%%%%%%%%%%%%%%%%%%
\balance
%%%%%%%%%%%%%%%%%%%%%%%%%%%%%%%%%%%%%%%%%%%%%%%%%%%%%%%%%%%%%%%%%



\medskip 
\question~({\bf Plano amostral geral})    Seja 
${\cal U}=\{1,2,3\}$ uma população finita de tamanho $N=3$ e $\mathbf{Y} =\{2,4,6\}$  o vetor da característica populacional renda bruta (mensal em salários mínimos) familiar.  Suponha que o seguinte plano amostral é implementado 
$p(s_1)=p(\{1,2\})=\frac{1}{2},$ $p(s_2)=p(\{1,3\})=\frac{1}{4}$ e
$p(s_3)=p(\{2,3\})=\frac{1}{4}.$
\begin{description}
% \item[a)]~\textbf{[2pts]} Determine as probabilidades de inclusão de primeira ordem. 
\item[a)] Determine as probabilidades de inclusão de primeira e  segunda ordem. Determine se o plano amostral induzido pelo esquema de amostragem proposto é mensurável? 
% \item[c)]~\textbf{[2pts]} Determine as covariâncias dos indicadores de pertencimento à amostra  $$I_k=I_k(s)=\begin{cases} 1, & k \in s, \\ 0, & \text{caso contrário}.\end{cases}$$
\item[b)] Forneça a distribuição de probabilidades do estimador de Horvitz-Thompson $\hat{t}_\pi$ para o total populacional da renda bruta (mensal em salários mínimos) familiar.
\item[c)] Determine a variância do estimador de Horvitz-Thompson $\hat{t}_\pi$ para o total populacional da renda bruta (mensal em salários mínimos) familiar sob este plano amostral.
\end{description}

{\scriptsize
\noindent \textbf{Res: 1a.} Primeiro, vamos definir as probabilidades de inclusão de primeira ordem. A probabilidade de inclusão de primeira ordem \(\pi_i\) é a probabilidade de um elemento \(i\) ser incluído na amostra. 

Para \(i = 1\):
\[
\pi_1 = P(\{1, 2\}) + P(\{1, 3\}) = \frac{1}{2} + \frac{1}{4} = \frac{3}{4}
\]

Para \(i = 2\):
\[
\pi_2 = P(\{1, 2\}) + P(\{2, 3\}) = \frac{1}{2} + \frac{1}{4} = \frac{3}{4}
\]

Para \(i = 3\):
\[
\pi_3 = P(\{1, 3\}) + P(\{2, 3\}) = \frac{1}{4} + \frac{1}{4} = \frac{1}{2}
\] 
A probabilidade de inclusão de segunda ordem \(\pi_{ij}\) é a probabilidade de os elementos \(i\) e \(j\) serem incluídos simultaneamente na amostra.

Para \(\pi_{12}\):
\[
\pi_{12} = P(\{1, 2\}) = \frac{1}{2}
\]

Para \(\pi_{13}\):
\[
\pi_{13} = P(\{1, 3\}) = \frac{1}{4}
\]

Para \(\pi_{23}\):
\[
\pi_{23} = P(\{2, 3\}) = \frac{1}{4}
\]

Um plano amostral é mensurável se todas as probabilidades de inclusão de segunda ordem forem maiores que zero e se \(\pi_{ij} \leq \min(\pi_i, \pi_j)\). Neste caso, temos:

\[
\pi_{12} = \frac{1}{2} \quad \text{e} \quad \min(\pi_1, \pi_2) = \frac{3}{4}
\]
\[
\pi_{13} = \frac{1}{4} \quad \text{e} \quad \min(\pi_1, \pi_3) = \frac{1}{2}
\]
\[
\pi_{23} = \frac{1}{4} \quad \text{e} \quad \min(\pi_2, \pi_3) = \frac{1}{2}
\]

Como todas as probabilidades de inclusão de segunda ordem são menores ou iguais ao mínimo das probabilidades de inclusão de primeira ordem correspondentes, podemos concluir que o plano amostral é mensurável.
}

\medskip
{\scriptsize

\textbf{Res:  1b} O estimador de Horvitz-Thompson para o total populacional da renda bruta (mensal em salários mínimos) familiar é dado por:

\[
\hat{t}_\pi = \sum_{i \in s} \frac{y_i}{\pi_i}
\]

Onde:
- \(y_i\) é o valor da característica para a unidade \(i\),
- \(\pi_i\) é a probabilidade de inclusão da unidade \(i\),
- \(s\) é a amostra.

Dado o plano amostral e as probabilidades de inclusão calculadas anteriormente, vamos calcular \(\hat{t}_\pi\) para cada amostra possível:

Para a amostra \(s_1 = \{1, 2\}\):
\[
\hat{t}_\pi(s_1) = \frac{y_1}{\pi_1} + \frac{y_2}{\pi_2} = \frac{2}{\frac{3}{4}} + \frac{4}{\frac{3}{4}} = \frac{2 \times 4}{3} + \frac{4 \times 4}{3} = \frac{8}{3} + \frac{16}{3} = 8
\]

Para a amostra \(s_2 = \{1, 3\}\):
\[
\hat{t}_\pi(s_2) = \frac{y_1}{\pi_1} + \frac{y_3}{\pi_3} = \frac{2}{\frac{3}{4}} + \frac{6}{\frac{1}{2}} = \frac{8}{3} + 12 = \frac{8}{3} + \frac{36}{3} = \frac{44}{3} \approx 14.67
\]

Para a amostra \(s_3 = \{2, 3\}\):
\[
\hat{t}_\pi(s_3) = \frac{y_2}{\pi_2} + \frac{y_3}{\pi_3} = \frac{4}{\frac{3}{4}} + \frac{6}{\frac{1}{2}} = \frac{16}{3} + 12 = \frac{16}{3} + \frac{36}{3} = \frac{52}{3} \approx 17.33
\]

Agora, vamos determinar a distribuição de probabilidades do estimador \(\hat{t}_\pi\):

\[
P\left(\hat{t}_\pi = 8\right) = P(s_1) = \frac{1}{2}
\]
\[
P\left(\hat{t}_\pi = \frac{44}{3} \approx 14.67\right) = P(s_2) = \frac{1}{4}
\]
\[
P\left(\hat{t}_\pi = \frac{52}{3} \approx 17.33\right) = P(s_3) = \frac{1}{4}
\]

A distribuição de probabilidades do estimador \(\hat{t}_\pi\) é dada por:

\[
\hat{t}_\pi =
\begin{cases}
8, & \text{com probabilidade } \frac{1}{2}, \\
\frac{44}{3} \approx 14.67, & \text{com probabilidade } \frac{1}{4}, \\
\frac{52}{3} \approx 17.33, & \text{com probabilidade } \frac{1}{4}.
\end{cases}
\]
}

\medskip
{\scriptsize

\textbf{Res:  1c} Podemos dividir o cálculo da variância em partes menores para facilitar. Cálculo da Variância do Estimador de Horvitz-Thompson \(\hat{t}_\pi\) A variância do estimador de Horvitz-Thompson é dada por:
\[
\text{Var}(\hat{t}_\pi) = \sum_{i \in \mathcal{U}} \sum_{j \in \mathcal{U}} \left(\frac{y_i}{\pi_i}\right) \left(\frac{y_j}{\pi_j}\right) (\pi_{ij} - \pi_i \pi_j)
\]

Passo 1: Calcular os termos diagonais

Primeiro, calculamos os termos em que \(i = j\), ou seja, onde a soma envolve as mesmas unidades:
\[
\text{Var}(\hat{t}_\pi)_{\text{diagonal}} = \sum_{i \in \mathcal{U}} \left(\frac{y_i}{\pi_i}\right)^2 (1 - \pi_i)
\]

Os valores que temos são:

- \(\pi_1 = \frac{3}{4}\), \(\pi_2 = \frac{3}{4}\), \(\pi_3 = \frac{1}{2}\)

- \(y_1 = 2\), \(y_2 = 4\), \(y_3 = 6\)

Substituindo:
\[
\text{Var}(\hat{t}_\pi)_{\text{diagonal}} = \left(\frac{2}{\frac{3}{4}}\right)^2 \left(\frac{1}{4}\right) + \left(\frac{4}{\frac{3}{4}}\right)^2 \left(\frac{1}{4}\right) + \left(\frac{6}{\frac{1}{2}}\right)^2 \left(\frac{1}{2}\right)
\]

Agora, vamos calcular esses três termos separadamente:

1. Para \(i = 1\):
\[
\left(\frac{2}{\frac{3}{4}}\right)^2 \times \frac{1}{4} = \left(\frac{8}{3}\right)^2 \times \frac{1}{4} = \frac{64}{9} \times \frac{1}{4} = \frac{16}{9} \approx 1.78
\]

2. Para \(i = 2\):
\[
\left(\frac{4}{\frac{3}{4}}\right)^2 \times \frac{1}{4} = \left(\frac{16}{3}\right)^2 \times \frac{1}{4} = \frac{256}{9} \times \frac{1}{4} = \frac{64}{9} \approx 7.11
\]

3. Para \(i = 3\):
\[
\left(\frac{6}{\frac{1}{2}}\right)^2 \times \frac{1}{2} = 12^2 \times \frac{1}{2} = 144 \times \frac{1}{2} = 72
\]

Somando esses termos:
\[
\text{Var}(\hat{t}_\pi)_{\text{diagonal}} \approx 1.78 + 7.11 + 72 = 80.89
\]
Este é o valor parcial para a variância considerando apenas os termos diagonais. Vamos agora calcular os termos cruzados da variância do estimador de Horvitz-Thompson, onde \(i \neq j\).

\[
\text{Var}(\hat{t}_\pi)_{\text{cruzados}} = \sum_{i \neq j} \left(\frac{y_i}{\pi_i}\right) \left(\frac{y_j}{\pi_j}\right) (\pi_{ij} - \pi_i \pi_j)
\]
Aqui, precisamos calcular cada combinação de \(i\) e \(j\) onde \(i \neq j\). Lembrando que as probabilidades de segunda ordem \(\pi_{ij}\) foram calculadas anteriormente: \(\pi_{12} = \frac{1}{2}\), \(\pi_{13} = \frac{1}{4}\) e \(\pi_{23} = \frac{1}{4}\). O Termo para \(i = 1\) e \(j = 2\):
\[
\left(\frac{y_1}{\pi_1}\right) \left(\frac{y_2}{\pi_2}\right) (\pi_{12} - \pi_1 \pi_2) = \left(\frac{2}{\frac{3}{4}}\right) \left(\frac{4}{\frac{3}{4}}\right) \left(\frac{1}{2} - \frac{3}{4} \times \frac{3}{4}\right)
\]
\[
= \frac{8}{3} \times \frac{16}{3} \times \left(\frac{1}{2} - \frac{9}{16}\right) = \frac{128}{9} \times \left(\frac{8}{16}\right) = \frac{128}{9} \times \frac{1}{2} = \frac{64}{9} \approx 7.11
\]

O termo para \(i = 1\) e \(j = 3\): 
\[
\left(\frac{y_1}{\pi_1}\right) \left(\frac{y_3}{\pi_3}\right) (\pi_{13} - \pi_1 \pi_3) = \left(\frac{2}{\frac{3}{4}}\right) \left(\frac{6}{\frac{1}{2}}\right) \left(\frac{1}{4} - \frac{3}{4} \times \frac{1}{2}\right)
\] 
\[
= \frac{8}{3} \times 12 \times \left(\frac{1}{4} - \frac{3}{8}\right) = \frac{96}{1} \times \left(\frac{2}{8} - \frac{3}{8}\right) = \frac{96}{1} \times \left(-\frac{1}{8}\right) = -12
\]

O termo para \(i = 2\) e \(j = 3\): 
\[
\left(\frac{y_2}{\pi_2}\right) \left(\frac{y_3}{\pi_3}\right) (\pi_{23} - \pi_2 \pi_3) = \left(\frac{4}{\frac{3}{4}}\right) \left(\frac{6}{\frac{1}{2}}\right) \left(\frac{1}{4} - \frac{3}{4} \times \frac{1}{2}\right)
\] 
\[
= \frac{16}{3} \times 12 \times \left(\frac{1}{4} - \frac{3}{8}\right) = \frac{192}{1} \times \left(\frac{2}{8} - \frac{3}{8}\right) = \frac{192}{1} \times \left(-\frac{1}{8}\right) = -24
\] 
Agora, somamos os termos cruzados: 
\[
\text{Var}(\hat{t}_\pi)_{\text{cruzados}} = 7.11 - 12 - 24 = -28.89
\] 
Finalmente, somamos os termos diagonais e cruzados para obter a variância total do estimador de Horvitz-Thompson:
\[
\text{Var}(\hat{t}_\pi) \approx 80.89 - 28.89 = 52
\] 
Este é o valor da variância do estimador de Horvitz-Thompson \(\hat{t}_\pi\) para o total populacional da renda bruta sob o plano amostral fornecido.

}



\medskip 
\question~~({\bf Plano amostral AASs}) Uma amostra aleatória simples e sem substituição de 56 pessoas foi selecionada de uma população de 1000 trabalhadores da empresa LInCaTech. Foram coletadas a informações sobre a Renda mensal em miles de reais (Renda) e o sexo do trabalhador. Com as informações da tabela \ref{ID:Renda} estime:

\begin{itemize}
    \item[a)] A renda média dos trabalhadores. Estabeleça um intervalo de 95\% para a renda média.
\item[b)] A renda total dos trabalhadores. Estabeleça um intervalo de 95\% para a renda total. 
\end{itemize}
{\small \begin{table}[!htp]
\centering
\begin{tabular}{rlr|rlr}
  \hline
ID & Sexo & Renda & ID & Sexo & Renda \\ 
  \hline
    1 & Fem & 2094.90 &   15 & Mas & 2939.21 \\
    2 & Mas & 2386.14 &   16 & Mas & 1722.62 \\
    3 & Mas & 1562.82 &   17 & Fem & 2739.79 \\
    4 & Mas & 1781.41 &   18 & Mas & 1821.61 \\
    5 & Mas & 1603.14 &   19 & Mas & 1742.40 \\
    6 & Mas &  479.73 &   20 & Mas & 1845.22 \\
    7 & Mas & 2196.85 &   21 & Mas & 1916.60 \\
    8 & Mas & 2365.21 &   22 & Mas & 1329.28 \\
    9 & Mas & 2016.32 &   23 & Mas & 2143.04 \\
   10 & Mas & 1322.23 &   24 & Mas & 2618.97 \\
   11 & Fem & 2589.08 &   25 & Fem & 1399.85 \\
   12 & Fem & 2896.28 &   26 & Mas & 1610.41 \\
   13 & Mas & 1370.55 &   27 & Mas & 2300.84 \\
   14 & Mas &  975.94 &   28 & Mas & 1192.90 \\ \hline   
   29 & Mas & 2715.62 &   43 & Mas & 1802.16 \\
   30 & Mas & 2042.58 &   44 & Mas & 2444.68 \\
   31 & Mas & 2235.73 &   45 & Mas & 2644.75 \\
   32 & Mas & 2223.33 &   46 & Fem & 1431.53 \\
   33 & Mas & 2618.16 &   47 & Mas & 1094.02 \\
   34 & Mas & 2206.57 &   48 & Mas & 1548.96 \\
   35 & Mas & 2432.08 &   49 & Mas & 2410.51 \\
   36 & Mas & 1340.94 &   50 & Mas & 2286.37 \\
   37 & Mas & 2321.37 &   51 & Mas & 1589.07 \\
   38 & Mas & 1922.91 &   52 & Fem & 1646.21 \\
   39 & Mas & 2520.83 &   53 & Mas & 3358.63 \\
   40 & Mas & 2063.78 &   54 & Mas & 1369.42 \\
   41 & Mas & 2335.66 &   55 & Mas & 2047.63 \\
   42 & Mas & 2357.94 &   56 & Mas & 1719.34 \\
   \hline
\end{tabular}
\caption{Tabela de Informações dos empregados na amostra} 
\label{ID:Renda}
\end{table}
}

{\scriptsize

{\bf Res: 2a} Cálculo da renda média e intervalo de confiança de 95\% para a renda média. O Cálculo da renda média amostral (\(\bar{X}\)) é 
   \[
   \bar{X} = \frac{1}{n} \sum_{i=1}^{n} X_i,
   \]
   onde \(n\) é o número de observações e \(X_i\) é a renda de cada trabalhador.   O desvio padrão amostral é calculado como:
   \[
   s = \sqrt{\frac{1}{n-1} \sum_{i=1}^{n} (X_i - \bar{X})^2}
   \]
   O intervalo de confiança de 95\% para a média amostral é dado por:

   \[
   \bar{X} \pm t_{\alpha/2, n-1} \cdot \frac{s}{\sqrt{n}}
   \]

   onde \(t_{\alpha/2, n-1}\) é o valor crítico da distribuição $t$ de Student com \(n-1\) graus de liberdade para um intervalo de confiança de 95%.

Vamos calcular esses valores passo a passo:

- Número de observações (\(n\)) = 56

- Rendas (\(X_i\)) fornecidas na tabela


1. Renda média amostral (\(\bar{X}\)):
\[
   \bar{X} = \frac{1}{56} \sum_{i=1}^{56} X_i
   \]

   % Primeiro, somamos todas as rendas:

   % \[
   % \text{Soma das rendas} = 2094.90 + 2386.14 + \cdots + 1719.34 = 110,538.54
   % \]
   Então,
   \[
   \bar{X} = \frac{110,538.54}{56} = 1,970.67
   \]

2. Desvio padrão amostral (\(s\)): Calculamos a variância amostral e depois tiramos a raiz quadrada:
\[
   s^2 = \frac{1}{55} \left[ (2094.90 - 1970.67)^2 + \cdots + (1719.34 - 1970.67)^2 \right] 
   \] Portanto,
   \[
   s = \sqrt{1,618,856.16} = 1,272.35
   \]

3. Intervalo de confiança de 95\%: Utilizando o valor crítico da distribuição t com 55 graus de liberdade (\(t_{0.025, 55} \approx 2.004\)): 
\[
   \text{Margem de erro} = 2.004 \cdot \frac{1272.35}{\sqrt{56}} = 2.004 \cdot 170.84 = 341.44
   \]
   O intervalo de confiança é: 
   \[
   \bar{X} \pm \text{Margem de erro} = 1970.67 \pm 341.44
   \] 
   \[
   \text{Intervalo de confiança de 95\%} = [1,629.23, 2,312.11]
   \]

{\bf Res: 2b} Cálculo da renda total e intervalo de confiança de 95\% para a renda total

Para calcular a renda total e seu intervalo de confiança, usamos: 
Renda total amostral (\(\hat{T} = n \cdot \bar{X}\)): 
\[
   \hat{T} = 56 \cdot 1,970.67 = 110,538.54
   \] 
   Utilizamos a fórmula para o intervalo de confiança da renda total:
   \[
   \hat{T} \pm t_{\alpha/2, n-1} \cdot \frac{s \cdot \sqrt{n}}{\sqrt{N - n}}
   \]
   O erro padrão da renda total é calculado como: 
   \[
   \text{Erro padrão} = t_{\alpha/2, n-1} \cdot \frac{s \cdot \sqrt{56}}{\sqrt{1000 - 56}}  = 564.46
   \] 
   O intervalo de confiança para a renda total é: 
   \[
   \hat{T} \pm \text{Erro padrão} = 110,538.54 \pm 564.46
   \]
   \[
   \text{Intervalo de confiança de 95\% para a renda total} = [109,974.08, 111,102.99]
   \]
}

\medskip

\question~({\bf Tamanho de amostra}) 
A variância da média amostral num plano amostral aleatório simples sem substituição (AASs) está dada por
$$
{\rm Var}_{_{\rm AASs}}(\bar{y}) = \Big( 1-\frac{n}{N}\Big)\frac{S^2_{Y_U}}{n},
$$
em que $S^2_{Y_U}$ é a variância populacional da variável de interesse, $n$ o tamanho de amostra e $N$ o tamanho populacional. Demonstre que o tamanho de amostra $n$ que garante uma variância máxima, estabelecida de antemão, ${\rm Var}_{_{\rm AASs}}(\bar{y})$ é
$$
n = \frac{
            \frac{S^2_{Y_U}}{{\rm Var}_{_{\rm AASs}}(\bar{y})}}{1+\frac{S^2_{Y_U}}{{\rm Var}_{_{\rm AASs}}}}.
$$

\medskip
{\scriptsize  {\bf Res: 4}
A variância da média amostral em um plano amostral aleatório simples sem substituição (AASs) é dada por
\[
{\rm Var}_{_{\rm AASs}}(\bar{y}) = \left( 1 - \frac{n}{N} \right) \frac{S^2_{Y_U}}{n},
\]
onde \( S^2_{Y_U} \) é a variância populacional, \( n \) é o tamanho da amostra, e \( N \) é o tamanho populacional. Queremos encontrar o tamanho de amostra \( n \) que garante uma variância máxima estabelecida de antemão, denotada por \( \sigma^2_{\bar{y}} \). Substituindo \( \sigma^2_{\bar{y}} \) na fórmula, temos:
\[
\sigma^2_{\bar{y}} = \left( 1 - \frac{n}{N} \right) \frac{S^2_{Y_U}}{n}.
\] Multiplicando ambos os lados por \( n \):
\[
n \sigma^2_{\bar{y}} = \left( 1 - \frac{n}{N} \right) S^2_{Y_U}.
\] Distribuindo \( S^2_{Y_U} \) no lado direito:
\[
n \sigma^2_{\bar{y}} = S^2_{Y_U} - \frac{n S^2_{Y_U}}{N}.
\] O que implica,  \[
n \sigma^2_{\bar{y}} + \frac{n S^2_{Y_U}}{N} = S^2_{Y_U}.
\] logo,  resolvendo para \( n \):
\[
n = \frac{S^2_{Y_U}}{\sigma^2_{\bar{y}} + \frac{S^2_{Y_U}}{N}}.
\] e  simplificando:
\[
n = \frac{\frac{S^2_{Y_U}}{\sigma^2_{\bar{y}}}}{1 + \frac{S^2_{Y_U}}{\sigma^2_{\bar{y}} N}}.
\]
}

\medskip 
\question~({\bf Amostragem Bernoulli}) Seja $s$ uma amostra obtida de um plano amostral de tipo Bernoulli com probabilidades de inclusão $\pi_k = \pi$ para todo $k \in U$ (população). Seja $n_s$ o tamanho de amostra da amostra $s.$ Mostre que a probabilidade condicional de se obter $s$ dado $n_s$ é a mesma que a probabilidade obtida por uma amostragem aleatória simples sem substituição de tamanho fixado $n_s$ de $N$ (Tamanho da população). 

\medskip
{\scriptsize  {\bf Res: 4}
Considere um plano amostral Bernoulli com probabilidades de inclusão \( \pi_k = \pi \) para todo \( k \in U \). Se \( s \) é uma amostra de tamanho \( n_s \), a probabilidade de observar essa amostra específica \( s \) é dada por: 
\[
P(s \mid n_s) = \pi^{n_s} (1 - \pi)^{N - n_s}.
\] Onde \( \pi^{n_s} \) é a probabilidade de incluir exatamente essas \( n_s \) unidades na amostra, e \( (1 - \pi)^{N - n_s} \) é a probabilidade de não incluir as restantes \( N - n_s \) unidades.

Por outro lado, na amostragem aleatória simples sem substituição (AASs), a probabilidade de obter uma amostra específica \( s \) de tamanho \( n_s \) é dada por:
\[
P(s) = \frac{1}{\binom{N}{n_s}},
\] 
onde \( \binom{N}{n_s} \) é o número total de combinações possíveis de tamanho \( n_s \) a partir da população de tamanho \( N \). 
Para amostras de tamanho fixo \( n_s \), a probabilidade de se obter uma amostra \( s \) é a mesma tanto para um plano amostral Bernoulli quanto para a amostragem aleatória simples sem substituição: 
\[
P(s \mid n_s) = \frac{1}{\binom{N}{n_s}} = P(s).
\] Isso ocorre porque, para amostras de tamanho fixo \( n_s \), todas as amostras possíveis têm a mesma probabilidade de serem selecionadas, seja em um plano amostral Bernoulli ou em uma amostragem aleatória simples sem substituição.
}


\medskip
\question~({\bf Amostragem Sistemática}) Suponha uma população de 9 elementos cujos valores para a característica de interesse sejam dados por
$\mathbf{Y} = \{23, 20, 24, 31, 24, 29, 25, 33, 21\}.$ Use a análise de variância (ANOVA) para calcular a variância do estimador de Horvitz-Thompson para o total populacional em um plano amostral sistemático com $a = 2$ grupos. Para esse caso específico, o plano amostral sistemático é mais eficiente do que o plano amostral aleatório simples sem reposição? Explique.


\medskip
{\scriptsize  {\bf Res: 5}
Temos a população:
\[
\mathbf{Y} = \{23, 20, 24, 31, 24, 29, 25, 33, 21\}
\] A média populacional é:

\[
\bar{Y} = \frac{1}{9} \sum_{i=1}^{9} Y_i = \frac{23 + \cdots  + 21}{9} = \frac{187}{9} = 20.77
\] 
A Variância da População 
\[
S^2 = \frac{1}{9-1} \sum_{i=1}^{9} (Y_i - \bar{Y})^2
\]

\[
S^2 = \frac{4.58 + 5.23 + \cdots + 0.64}{8} = 43.62
\] 
Plano Amostral Sistemático. Para um plano amostral sistemático com \(a = 2\) grupos, a variância do estimador é:
\[
\text{Var}(\hat{T}_{S}) = \frac{N - n}{n} \left( \frac{S^2_{B}}{a} + \frac{S^2_{W}}{n} \right)
\] em que \(N = 9\), \(n = 2\) e \(a = 2\). A Variância Entre os Grupos \(S^2_{B}\) é
\[
S^2_{B} = \frac{1}{a-1} \left[ (\bar{Y}_{1} - \bar{Y})^2 + (\bar{Y}_{2} - \bar{Y})^2 \right]
\] 
Com médias dos grupos:

- Grupo 1: \( \{23, 24, 24, 25, 21\} \) com média \( \bar{Y}_{1} = 23.4 \)

- Grupo 2: \( \{20, 31, 29, 33\} \) com média \( \bar{Y}_{2} = 28.25 \)

\[
S^2_{B} = \frac{1}{2-1} \left[ (23.4 - 20.77)^2 + (28.25 - 20.77)^2 \right]
\]
\[
S^2_{B} = \left[ 7.78^2 + 7.48^2 \right] = 60.5
\]
Por outro lado, a Variância Dentro dos Grupos \(S^2_{W}\):
\[
S^2_{W} = \frac{1}{N - a} \left[ \sum_{i=1}^{a} \sum_{j=1}^{n_i} (Y_{ij} - \bar{Y}_{i})^2 \right]
\]

Para o Grupo 1 e Grupo 2, calculamos:

- Grupo 1: \(S^2_{W1} = 6.2\)

- Grupo 2: \(S^2_{W2} = 40.7\)
\[
S^2_{W} = \frac{5 \cdot 6.2 + 4 \cdot 40.7}{9 - 2} = \frac{31 + 162.8}{7} = 27.8
\]

Assim, a Variância do Estimador Sistemático é
\[
\text{Var}(\hat{T}_{S}) = \frac{9 - 2}{2} \left( \frac{60.5}{2} + \frac{27.8}{2} \right) = \frac{7}{2} \cdot 44.15 = 154.52
\] 
Para o plano Amostral Aleatório Simples Sem Reposição 
\[
\text{Var}(\hat{T}_{S}) = \frac{N - n}{N - 1} \cdot \frac{S^2}{n}
\]
\[
\text{Var}(\hat{T}_{S}) = \frac{9 - 2}{9 - 1} \cdot \frac{43.62}{2} = \frac{7}{8} \cdot 21.81 = 19.21
\]

Quando comparamos as variâncias:

- Plano Sistemático: 154.52

- Plano Aleatório Simples: 19.21

Portanto, a variância do estimador de Horvitz-Thompson no plano amostral sistemático é significativamente maior do que no plano amostral aleatório simples sem reposição. Assim, o plano amostral aleatório simples sem reposição é mais eficiente neste caso específico.

\vspace{17cm}
}
\vfill\eject

\end{document}


